Overall, the experiment demonstrates that a simple optical setup,
supported by accessible Arduino-based instrumentation, can determine
refractive indices across a range of liquids.
The LDR-based detection system was able to distinguish between liquids
with similar refractive indices, demonstrating sensitivity to changes
on the order of $\Delta n \approx 0.03$.
Using the optimized quartz macro cuvette setup, water was calculated to
have a refractive index of $1.424 \pm 0.211$ (532~nm laser) with a
percent error of 6.8\% from the theoretical value of 1.333.
Ethanol measurements ranged from $1.220 \pm 0.077$ to
$1.533 \pm 0.197$ depending on wavelength and setup, yielding percent
errors between 10.4\% and 12.6\% from the theoretical value of 1.361.
Extra virgin olive oil was determined to be $1.175 \pm 0.140$ (532~nm)
with approximately 20\% underestimation from the theoretically accepted
range of 1.465--1.480.

While the quantitative accuracy was limited by systematic errors in
fringe detection and optical alignment, the experiment successfully
validated the fundamental principle of diffraction-based refractometry
and demonstrated the feasibility of low-cost instrumentation for
optical characterization.
The results highlight both the promise and the challenges of
Arduino-based precision optics measurements, suggesting that with
further refinement of the detection algorithm and mechanical stability,
this approach could provide a practical educational tool and
potentially a viable alternative for industrial applications requiring
moderate precision.

